\section{Introduction}
\label{sec:intro}

\begin{figure}[t]
\centering
\includegraphics[width=\linewidth]{figures/fig1.pdf}
\caption{Overview of TwinRouterBench. The benchmark provides a fast
static track for offline router development and a live dynamic track
for end-to-end validation. The static track covers 970 step-level rows
from 520 instances across five workloads, each with an
execution-verified target tier; the dynamic track runs routers on
SWE-bench Verified with realized API cost.}
\label{fig:overview}
\end{figure}

LLM routers select which model should handle each
call~\citep{feng2024graphrouter,stripelis2024tensoropera}. In production
multi-turn agents, each call carries a rich prefix---chat history, retrieved
passages, shell logs, tool outputs, and partial code edits---and a cheaper
choice at step $i$ may only fail downstream. Existing benchmarks evaluate
routing on isolated one-shot prompts; they do not expose these prefixes or
test whether a per-step decision breaks a later step.

The expensive unit in modern applications is rarely a one-shot query. Coding
agents solve GitHub issues through many rounds of exploration, patching, and
verification \citep{jimenez2023swebench}; tool systems maintain state over
serial and parallel calls \citep{patil2025bfcl}; retrieval and summarization
workloads carry long multi-turn context
\citep{katsis2025mtrag,zhong2021qmsum}. A cheap model may suffice for a file
lookup, while a later patch-writing step may require the strongest tier. A
useful step-level benchmark must cover these workloads, ground labels in
execution outcomes rather than model-based judgement alone, and support both
cheap deterministic offline scoring and end-to-end execution with realized
API cost as the definitive validation. We therefore maintain two distinct
tracks.

RouterBench, LLMRouterBench, and RouterArena evaluate query-level routing on
complete one-shot prompts
\citep{hu2024routerbench,li2026llmrouterbench,lu2025routerarena}. Such a
router can still be applied mid-trajectory by feeding the full conversation
history as a single query, but its training and evaluation data are one-shot
prompts, so neither tells us whether a cheap choice at step $i$ breaks a
later step. We expose the full prefix a router observes before each LLM call (chat
history, retrieved passages, tool outputs, shell traces, and code diffs)
as a first-class input, and we verify each label by running the task to
completion: if the trajectory still passes with the same number of steps
after downgrading, we infer the replaced step was handled correctly.
TRIM studies reasoning-step escalation in math
\citep{kapoor2026trim}; our setting is agentic and multi-turn.

Building such a benchmark faces two main challenges.
First, using a frontier model as the router is unreliable: in our
Claude Opus 4.6 diagnostic, a single-template LLM router predicts
\texttt{high} for only 7 of 147 verified-\texttt{high} SWE-bench steps and
fails every SWE trajectory.
Second, multi-turn trajectories create combinatorial label complexity: with
$n$ candidate models and $m$ steps, exhaustive verification requires $n^{m}$
runs even when the trajectory length is fixed, because the prefix at step
$i{+}1$ depends on the model chosen at step~$i$. Practical label
construction must therefore approximate this space with
affordable protocols.

We introduce \textbf{TwinRouterBench}, a step-level routing benchmark with a
fast static track for offline development and a live dynamic track for
end-to-end validation.

\paragraph{Contributions.}
\begin{itemize}[leftmargin=*, itemsep=1pt, topsep=2pt]
\item \textbf{A fast static track for realistic routing supervision.}
The static track covers the workloads where step-level routing pays off
most (long-context, multi-turn, tool-use, RAG, summarization, and
code-repair) and exposes the exact prefix a router observes before each LLM
call. Its
970 labels are execution-verified estimates under a fixed
downgrade-and-cascade protocol with tier-pool cascade verification,
mixed-prefix reconstruction, and manual audit on multi-step executable
workloads. Scoring is deterministic arithmetic over labels, trajectory
membership, and token costs: evaluation takes milliseconds and requires no
online evaluator-side LLM judge. The same data also supports training; a
logistic router trained on the static labels achieves the best cost--success
trade-off among the policies we evaluate.
\item \textbf{A live dynamic track for software-agent routing.}
The dynamic harness supports the full 500-case SWE-bench Verified suite; this
paper reports a 100-case held-out evaluation disjoint from the static SWE
supervision split. At each LLM request the router selects a concrete model
from a locked pool; scoring uses the official SWE-bench resolution predicate,
realized provider spend, cache accounting, and a flat unresolved-instance
penalty in the leaderboard bill (\S\ref{sec:dynamic-scoring}), with no
evaluator-side LLM judge.
\item \textbf{A two-track development and validation loop.}
The static track provides cheap, trainable supervision for rapid iteration.
The dynamic track tests whether the resulting policy survives live execution.
We report them side by side so that offline and end-to-end evidence remain
distinguishable.
\end{itemize}
